\section{第03章 ソフトウェア設計}
\begin{pointbox}{この資料の主張}
  ソフトウェア設計は、要求を単に「コードにする前の図」に変換する作業ではない。設計は、要求・制約・利害関係者の関心を、実装でき、評価でき、保守できる解決策へ変換する活動である。設計では、抽象化、関心事の分離、モジュール化、カプセル化、結合度、凝集度などの原則を使って複雑さを扱う。また、設計判断には根拠があり、その根拠を記録することで、後の変更や保守で同じ議論を繰り返さずに済む。
  \end{pointbox}

\begin{pointbox}{この資料の読み方}
  専門用語は原則として日本語で示し、必要に応じて英語を括弧内に併記する。英語名は、元資料やツール上の名称を確認したいときの手がかりとして残す。初学者が前提知識なしに読めるように、各節では「何のための概念か」「設計時に何を決めるのか」「実務では何を確認するのか」を重視する。文体は、だである調で統一する。
  \end{pointbox}
